\chapter{Performance Tuning}
\label{ch:performance}

This chapter details the performance configurations implemented in RGUKT Connect and recommends tuning strategies for future scaling.

\section{Implemented Performance Configurations}
\begin{itemize}
    \item \textbf{JPA Lazy Fetching:} Entity mappings use lazy fetching (\texttt{FetchType.LAZY}) for collections (such as `likedBy` in `Post` and `likedBy` in `Comment`). This prevents the system from loading related records until they are requested.
    \item \textbf{HikariCP Connection Pool:} Database connections are managed by a HikariCP pool. HikariCP reduces connection latency by keeping a pool of active database connections open, avoiding the overhead of establishing a new TCP connection for every request.
    \item \textbf{Binary Offloading to S3:} Profile pictures and post attachments are offloaded directly to Amazon S3. The database only stores the resulting URL strings, keeping database size small and reducing query latency.
    \item \textbf{Index-Optimized Relationships:} Core tables use auto-generated indexes on primary and foreign keys, speeding up join operations.
\end{itemize}

\section{Performance Tuning Recommendations}
\begin{enumerate}
    \item \textbf{Resolve N+1 Query Issues:} When loading the community feed, the application maps each `Post` entity to a `PostResponseDTO`. This maps user details, comment counts, and like states, which can trigger additional queries for each post in the list (the N+1 query problem). Using entity graphs or fetch joins (`JOIN FETCH`) in custom JPQL queries would allow the system to load this data in a single query.
    \item \textbf{Query Pagination:} Endpoints like `GET /api/posts` and `GET /api/jobs` load all records in a single request. As the platform grows, loading all records will degrade response times. The endpoints should use Spring Data's `Pageable` interface to paginate results (e.g., loading 20 posts at a time).
    \item \textbf{API Caching with Redis:} Read-heavy endpoints that change infrequently (such as the alumni directory directory search and job board listings) should be cached using Redis. This reduces database load and speeds up response times.
    \item \textbf{Frontend Asset Optimizations:} The React client should implement code splitting using `React.lazy` and `Suspense` to load route components on demand, reducing the size of the initial JavaScript bundle.
\end{enumerate}
